\documentclass[10pt,a4paper]{article}
\usepackage{amsmath,amssymb,amsthm,graphicx}
\usepackage{hyperref}
\usepackage{cleveref}
\usepackage{amssymb,amsmath,amsthm}
\newcommand{\Mod}[1]{\ (\mathrm{mod}\ #1)}
\newcommand\blfootnote[1]{%
  \begingroup
  \renewcommand\thefootnote{}\footnote{#1}%
  \addtocounter{footnote}{-1}%
  \endgroup
}

\linespread{1.1}



\DeclareMathOperator{\Tr}{Tr}
\newtheorem{theorem}{Theorem}[section]
\newtheorem{lemma}{Lemma}[section]
\newtheorem{proposition}{Proposition}[section]
\newtheorem{corollary}{Corollary}[section]
\newtheorem{assumption}{Assumption}
\newtheorem{fact}{Fact}[section]
\newtheorem{claim}{Claim}
\newtheorem{question}{Question}
\newtheorem{remark}{Remark}
\newtheorem{definition}{Definition}[section]
\title{The systematic scan contracts at a rate governed by the uniform gap parameter}
\author{Jason Gaitonde, Elchanan Mossel}
\date{arxiv:2410.11136v2; CC BY 4.0}
\begin{document}
\maketitle
\noindent\textit{Editorial note.} The counted statement is the author's labelled environment at source0.tex lines 535-541 and its proof environment opens at line 542; the pairing is by adjacency and the retained proof range is the whole environment.\par\medskip
\noindent\textit{Editorial note.} This source is newly approved for this line and its preamble is built from the author's own package line, macro declarations and theorem-environment family (source lines 4-70); the drawing, algorithm and font packages of the source are not retained because the retained material does not use them.\par\medskip
\noindent\textit{Editorial note.} Metadata provenance: the pinned source carries a licence observation file (CC BY 4.0, status supported, authors recorded there) but no separate source record file; per the project rule the missing metadata falls back to the licence observation, and the card records that here.\par\medskip
\noindent\textit{Editorial note.} The source carries no label-number dump and no counter values are imposed, so the printed numbering is the wrapper's default; identity is the author label.\par\medskip
\noindent\textit{Editorial note.} r002 note: r001 failed because the retained context window had been taken from a planner proposal and began inside the preceding proof, so the wrapper emitted an unmatched proof environment; r002 replaces the windows by the true enclosing environments of the cited labels (thm:gen\_gap at lines 516-531, cor:spectral\_scan at lines 535-541, thm:mixing at lines 126-133) and changes nothing else.\par\medskip
\noindent\textit{Editorial scope.} One original result of the pinned author manuscript, counted once: the contraction rate of the systematic scan under the general spectral gap assumptions (author label cor:spectral\_scan).\par\medskip
\noindent\textit{Allowed standard prerequisites (editorial).}\begin{itemize}
\item Assumptions of the general gap theorem. A family of coordinate transition matrices whose parameters satisfy the spectral independence bounds of the general theorem.. They are inherited by the corollary.
\item Uniform gap parameter. The parameter delta entering the general gap theorem bounds the spectral gap of every coordinate step.. It is the quantity appearing in the contraction rate.
\item Operator norm of the product. The norm of the product of the coordinate matrices measured on mean-zero functions.. It is the quantity bounded.
\end{itemize}\medskip
\begin{theorem}
\label{thm:gen_gap}
        Let $\mathcal{H}$ be a Hilbert space and suppose $\mathsf{P}_1,\ldots, \mathsf{P}_n$ are orthogonal projections with respect to the inner product of $\mathcal{H}$ satisfying
        \begin{equation}
    \left\|\frac{1}{n}\sum_{i=1}^n \mathsf{P}_i\right\|_{\mathsf{op}}\leq 1-\delta.
\end{equation} 
Let $\mathsf{P}_{i_L}\ldots \mathsf{P}_{i_1}$ be any product of the orthogonal projection operators satisfying the following condition. For each $j\in [n]$, suppose there exists a minimal index $k_j\leq L$ satisfying $i_{k_j}=j$. We further assume these indices satisfy\footnote{In words, the indices $i_1,\ldots,i_L$ contain $[n]$ in ascending order of first appearance and are a minimal subsequence with this property.}:
\begin{equation}
\label{eq:order}
        1=k_1\leq k_2\leq\ldots\leq k_n=L.
\end{equation}
Then, it holds that
\begin{equation*}
    \|\mathsf{P}_{i_L}\ldots \mathsf{P}_{i_1}\|_{\mathsf{op}}^2\leq 1-\frac{n\delta}{8\sum_{j=1}^n k_j}.
\end{equation*}
\end{theorem}

\begin{corollary}
\label{cor:spectral_scan}
    Under the assumptions of \Cref{thm:gen_gap}, it holds that 
    \begin{equation*}
        \|\mathsf{P}_n\ldots \mathsf{P}_1\|_{\mathsf{op}}\leq 1-\frac{\delta}{8(n+1)}.
    \end{equation*}
\end{corollary}

\begin{proof}
    For the systematic scan, it is clear that the assumptions of \Cref{thm:gen_gap} are satisfied since we have $k_j=j$ for each $j\leq n$. Therefore, \Cref{thm:gen_gap} implies that
    \begin{equation*}
        \|\mathsf{P}_n\ldots \mathsf{P}_1\|_{\mathsf{op}}^2\leq 1-\frac{n\delta}{8\sum_{j=1}^n j}=1-\frac{\delta}{4(n+1)}.
    \end{equation*}
    The desired conclusion then follows from Bernoulli's inequality.
\end{proof}

\end{document}
